\section{Evaluations}

\subsection{Datasets}
We used the two most popular ULVU datasets in the paper: EgoLifeQA~\cite{yang_egolife_2026} and MM-Lifelong~\cite{chen_multimodal_2026}. EgoLifeQA uses multiple-choice questions (MCQ) while MM-Lifelong uses open QA. EgoLifeQA contains one week of recording in the same setting while MM-Lifelong has 3 different settings (day, week, and month), each with its own video. The MM-Lifelong week subset has the same video as EgoLifeQA. We argue that MCQ does not reflect realistic usage since most questions asked by users in the real world are open. Additionally, agents can brute-force verify each option to artificially boost accuracy. We therefore reformulated EgoLifeQA into open QA questions (Open-EgoLifeQA) similar to MM-Lifelong and evaluated both using LLM-as-judge for semantic comparison between the ground-truth and predicted answers. To compare with prior work, we used the Jack view from EgoLifeQA. We evaluated our untrained VideoSQL on both Open-EgoLifeQA and the MM-Lifelong validation set.

To train the VideoSQL Skill, we need to select an independent training set. Ego-R1 Data~\cite{tian_egor1_2025} is a training set built on top of EgoLifeQA and the authors of Ego-R1 trained their method on this dataset and then tested on EgoLifeQA. MM-Lifelong also has a training split as well. However, these training sets overlap with their corresponding validation/testing sets at the video level. While no questions are leaked, we argue that sharing videos could still inflate performance, or at least fail to test the model's ability for cross-domain generalization. To address this, we trained our Skill on EgoLifeQA and tested it on the more challenging MM-Lifelong (for context, MAGIC-Video reached 67.6 on EgoLifeQA but only 24.5 on MM-Lifelong).


\subsection{Baselines}
We re-ran the MAGIC-Video \cite{li_bridging_2026} as our training-free baseline as we share the same indices. Both MAGIC-Video and our method used Qwen-3.8-27B~\cite{qwen38} as the policy model and Qwen3-VL-Embedding-8B \cite{li_qwen3vlembedding_2026} as embedding model. We also compared with a couple other baselines directly with their reported numbers. We also selected Trace2Skill \cite{ni_trace2skill_2026} and SkillOpt \cite{yang_skillopt_2026} for our skill-trained variant, both built on top of the same VideoSQL harness.

